\chapter{Foundations and Interfaces}

This chapter fixes three interfaces used by the surgery proof: comparison for
upper forward derivatives, the local geometry of neck surgery, and the
topology of regions covered by canonical necks and caps.

\section{Forward-Difference Comparison}

For a continuous function $f\colon[a,b]\to\mathbb R$, write
$D^+f(t)\le c$ if
\[
  \limsup_{h\downarrow0}\frac{f(t+h)-f(t)}h\le c,
  \qquad t\in[a,b).
\]

\begin{lemma}[Right-Dini mean-value lemma]
\label{lem:forward-dini-mean-value}
\dcref{MT:ch2:2.22}
\notready
Let $q\colon[c,d]\to\mathbb R$ be continuous.  If $q(d)>q(c)$, then
there is $s\in[c,d)$ such that $D^+q(s)>0$.
\end{lemma}

\begin{proof}
Choose $0<\eta<(q(d)-q(c))/(d-c)$ and let $s$ minimize
$q(t)-\eta(t-c)$ on $[c,d]$.  Since
\[
 q(d)-\eta(d-c)>q(c),
\]
the point $d$ is not a minimizer, so $s<d$.  For all sufficiently small
$h>0$, minimality gives $q(s+h)-q(s)\ge\eta h$.  Hence
$D^+q(s)\ge\eta>0$.
\end{proof}

\begin{theorem}[Forward-Dini comparison principle]
\label{thm:forward-dini-comparison}
\dcref{MT:ch2:2.22}
\notready
Let $f\colon[a,b]\to\mathbb R$ be continuous, let
$\psi\colon[a,b]\times\mathbb R\to\mathbb R$ be of class $C^1$, and
suppose
\[
  D^+f(t)\le \psi(t,f(t))\qquad(a\le t<b).
\]
If $G\in C^1([a,b])$ satisfies $G'(t)=\psi(t,G(t))$ and
$f(a)\le G(a)$, then $f(t)\le G(t)$ for every $t\in[a,b]$.
\end{theorem}

\begin{proof}
\uses{lem:forward-dini-mean-value}
The ranges of $f$ and $G$ lie in a compact interval $J$.  The derivative
$\partial_2\psi$ is bounded on $[a,b]\times J$, so there is $L\ge0$ such
that
\[
 |\psi(t,u)-\psi(t,v)|\le L|u-v|
 \quad(t\in[a,b],\ u,v\in J).
\]
Set $u=f-G$.  Whenever $u(t)\ge0$, subtraction of the right derivative of
$G$ from the upper forward derivative of $f$ gives
\[
 D^+u(t)\le \psi(t,f(t))-\psi(t,G(t))\le Lu(t).
\]
Consequently $v(t)=e^{-L(t-a)}u(t)$ satisfies $D^+v(t)\le0$ whenever
$v(t)\ge0$.

Suppose that $v(d)>0$ for some $d$.  Since $v(a)\le0$, continuity supplies
$c<d$ with $v(c)=0$ and $v>0$ on $(c,d]$.  By
\ref{lem:forward-dini-mean-value}, some $s\in[c,d)$ satisfies
$D^+v(s)>0$.  This contradicts $v(s)\ge0$ and the inequality above.
Thus $v\le0$, and hence $f\le G$.
\end{proof}

\begin{lemma}[Uniform transport of time-slice maxima]
\label{lem:uniform-time-slice-transport}
\dcref{MT:ch2:2.23}
\notready
Let $X$ be a set with functions $\tau,F,V\colon X\to\mathbb R$, let
$Z\subset X$, and let $m\colon[a,b]\to\mathbb R$.  Suppose that for every
$t\in[a,b]$ there is $x\in Z$ with $\tau(x)=t$ and $F(x)=m(t)$, and that
$F(y)\le m(\tau(y))$ whenever $\tau(y)\in[a,b]$.  Suppose also that there
are $\rho>0$ and transports $\Phi_h(x)$, defined for $x\in Z$,
$|h|<\rho$, and $\tau(x)+h\in[a,b]$, such that
\[
 \tau(\Phi_h(x))=\tau(x)+h,
 \qquad
 \sup_{x\in Z}|F(\Phi_h(x))-F(x)|\longrightarrow0
 \quad(h\to0).
\]
Assume moreover that, for every $\eta>0$, there is $h_0>0$ such that
\[
 \frac{F(x)-F(\Phi_{-h}(x))}{h}\le V(x)+\eta
\]
whenever $x\in Z$, $0<h<h_0$, and the backward transport is defined.
If $\psi\colon[a,b]\times\mathbb R\to\mathbb R$ is continuous and
\[
 V(x)\le\psi(\tau(x),F(x))\qquad(x\in Z),
\]
then $m$ is continuous and $D^+m(t)\le\psi(t,m(t))$ for every $t<b$.
\end{lemma}

\begin{proof}
For nearby $s,t\in[a,b]$, transport a maximizer $x_s\in Z$ from time $s$
to time $t$.  Since its transported value is at most $m(t)$,
\[
 m(s)-m(t)\le F(x_s)-F(\Phi_{t-s}(x_s)).
\]
Interchanging $s$ and $t$ and using the uniform value control proves the
continuity of $m$.

Fix $t<b$ and choose a maximizer $x_h\in Z$ at time $t+h$.  Backward
transport gives
\[
 \frac{m(t+h)-m(t)}h
 \le \frac{F(x_h)-F(\Phi_{-h}(x_h))}{h}.
\]
For every $\eta>0$, the right-hand side is at most
$V(x_h)+\eta$, and therefore at most
$\psi(t+h,m(t+h))+\eta$, for all sufficiently small $h>0$.
Letting $h\downarrow0$ and then $\eta\downarrow0$ proves the asserted
upper forward-derivative bound.
\end{proof}

\begin{theorem}[Time-slice maximum comparison]
\label{thm:time-slice-maximum-comparison}
\dcref{MT:ch2:2.23}
\notready
Let $M$ be a smooth manifold, let $\chi$ be a smooth vector field, and let
$\mathbf t\colon M\to[a,b]$ be smooth with $\chi(\mathbf t)=1$.  Let
$F\colon M\to\mathbb R$ be smooth, suppose that $F$ attains a maximum on
every level set of $\mathbf t$, and suppose that the set $\mathcal Z$ of all
levelwise maximizers is compact.  If
$\psi\in C^1([a,b]\times\mathbb R)$ and
\[
 \chi(F)(x)\le\psi(\mathbf t(x),F(x))
 \qquad(x\in\mathcal Z),
\]
then
\[
 F_{\max}(t)=\max_{\mathbf t^{-1}(t)}F
\]
is continuous and satisfies
$D^+F_{\max}(t)\le\psi(t,F_{\max}(t))$ for $t<b$.  Hence, if
$G'=\psi(t,G)$ and $F_{\max}(a)\le G(a)$, then
$F_{\max}(t)\le G(t)$ throughout $[a,b]$.
\end{theorem}

\begin{proof}
\uses{lem:uniform-time-slice-transport,thm:forward-dini-comparison}
Let $\varphi_h$ be the local flow of $\chi$.  Compactness of $\mathcal Z$
gives a common lifetime $\rho>0$ on a neighborhood of $\mathcal Z$, and
\[
 \mathbf t(\varphi_h(x))=\mathbf t(x)+h,
 \qquad
 \frac{d}{dh}F(\varphi_h(x))=\chi(F)(\varphi_h(x)).
\]
On a compact flow tube around $\mathcal Z$, smoothness gives both uniform
continuity of $F(\varphi_h(x))$ and the uniform expansion
\[
 \frac{F(x)-F(\varphi_{-h}(x))}{h}\longrightarrow\chi(F)(x).
\]
Thus \ref{lem:uniform-time-slice-transport}, with
$Z=\mathcal Z$, $V=\chi(F)$, and $m=F_{\max}$, proves continuity and the
forward-difference inequality.  The comparison with $G$ follows from
\ref{thm:forward-dini-comparison}.
\end{proof}

\section{Controlled-Surgery Inputs}

\subsection{Necks and singular-time extension}

\begin{definition}[Epsilon-closeness, neck, and scale]
\label{def:epsilon-neck-and-scale}
\dcref{MT:ch2:2.16,MT:ch2:2.18}
\notready
Fix $0<\epsilon<1/2$ and put $k=\lfloor\epsilon^{-1}\rfloor$.  A metric
$g$ is \emph{$\epsilon$-close to $g_0$ in $C^k$} on a common domain if
the $g_0$-norms of $g-g_0$ and its first $k$ covariant derivatives are
uniformly less than $\epsilon$.

An \emph{$\epsilon$-neck} centered at $x$ in a Riemannian three-manifold
$(M,g)$ is the image $N$ of a diffeomorphism
\[
 \varphi\colon S^2\times(-\epsilon^{-1},\epsilon^{-1})\longrightarrow N
\]
with $x\in\varphi(S^2\times\{0\})$, such that
$R(x)\varphi^*g$ is $\epsilon$-close in
$C^{\lfloor1/\epsilon\rfloor}$ to $h_0+ds^2$, where $h_0$ is the round
metric of scalar curvature $1$ on $S^2$.  The slices are the neck spheres,
the slice at $s=0$ is the central sphere, and
\[
 s_N=\operatorname{pr}_2\circ\varphi^{-1},
 \qquad r_N=R(x)^{-1/2}
\]
are the axial coordinate and the scale.  Replacing $s_N$ by $-s_N$
reverses the neck.  Its positive and negative ends are respectively
\[
 E_N^+=\{s_N>0.9\epsilon^{-1}\},
 \qquad
 E_N^-=\{s_N<-0.9\epsilon^{-1}\}.
\]
\end{definition}

\begin{definition}[Maximal extension of a generalized Ricci flow]
\label{def:maximal-ricci-flow-extension}
\dcref{MT:ch11:11.22}
\notready
Here a \emph{generalized Ricci flow} consists of a manifold with a time
function $\mathbf t$, a time vector field $\chi$ satisfying
$\chi(\mathbf t)=1$, and metrics on $\ker(d\mathbf t)$ which evolve by
the Ricci-flow equation along $\chi$.  Let $(\mathcal M,G)$ be such a
flow on a terminal interval
$[T^-,T)$, identify this part of the flow with $M\times[T^-,T)$, and set
\[
 \Omega=\{x\in M:\liminf_{t\uparrow T}R(x,t)<\infty\}.
\]
Suppose that $\Omega$ is open and that $g(t)|_\Omega$ converges smoothly
on compact subsets to a metric $g(T)$.  The \emph{maximal extension to
$T$} is
\[
 \widehat{\mathcal M}
 =\mathcal M\cup_{\Omega\times[T^-,T)}
   (\Omega\times[T^-,T]).
\]
The time function, time vector field, and horizontal metrics agree on the
overlap and hence glue to a generalized Ricci flow
$(\widehat{\mathcal M},\widehat G)$.  Its terminal time-slice is
$(\Omega,g(T))$, which need not be complete.
\end{definition}

\subsection{Local surgery data and metric}

\begin{definition}[Standard surgery-cap model]
\label{def:standard-surgery-cap-model}
\dcref{MT:ch12:12.2,MT:ch12:12.3}
\uses{def:epsilon-neck-and-scale}
\notready
Fix a complete rotationally symmetric metric $g_0$ on $\mathbb R^3$ with
nonnegative sectional curvature and tip $p_0$.  Fix constants
$A_0>r_0>0$ such that $g_0$ has constant sectional curvature $1/4$ on
$B_{g_0}(p_0,r_0)$ and there is an isometry
\[
 \psi\colon(S^2\times(-\infty,4],h_0+ds^2)
 \longrightarrow(\mathbb R^3\setminus B_{g_0}(p_0,A_0),g_0).
\]
Extend the cylinder coordinate to a continuous radial function by
\[
 s_1(y)=A_0+4-d_{g_0}(p_0,y).
\]
It is smooth away from $p_0$ and has round two-sphere level sets.
\end{definition}

\begin{definition}[Local neck-surgery data]
\label{def:local-neck-surgery-data}
\dcref{MT:ch13:13.1,MT:ch13:13.2}
\uses{def:epsilon-neck-and-scale,def:standard-surgery-cap-model}
\notready
Let $(N,g)$ be a $\delta$-neck centered at $x_0$, let
$Q=R_g(x_0)$, and let
\[
 \rho\colon S^2\times(-\delta^{-1},\delta^{-1})\longrightarrow N
\]
be its neck chart.  Assume $\delta^{-1}>A_0+4$.  Form $\mathcal S$ by
gluing $S^2\times(-\delta^{-1},4)$ to
$B_{g_0}(p_0,A_0+4)$ through the cylindrical isometry on
$S^2\times(0,4)$.  The two axial coordinates glue to
\[
 s\colon\mathcal S\longrightarrow(-\delta^{-1},A_0+4].
\]
Fix smooth cutoffs $\alpha\colon[1,2]\to[0,1]$ and
$\beta\colon[A_0+4-r_0,A_0+4]\to[0,1]$, equal to $1$ near their left
endpoints and $0$ near their right endpoints.  Put
\[
 \eta=1-2\delta,
 \qquad
 f(s)=\begin{cases}
 0,&s\le0,\\
 C_0\delta e^{-q/s},&s>0.
 \end{cases}
\]
\end{definition}

\begin{definition}[Interpolating surgery metric]
\label{def:interpolating-surgery-metric}
\dcref{MT:ch13:13.2}
\uses{def:local-neck-surgery-data}
\notready
Set $A_r=A_0+4-r_0$ and $A'=A_0+4$.  On $\mathcal S$, define
\[
\widehat g=
\begin{cases}
 e^{-2f}Q\rho^*g,&s\le1,\\
 e^{-2f}\bigl(\alpha Q\rho^*g+(1-\alpha)\eta g_0\bigr),
     &1\le s\le2,\\
 e^{-2f}\eta g_0,&2\le s\le A_r,\\
 \bigl(\beta e^{-2f}+(1-\beta)e^{-2f(A')}\bigr)\eta g_0,
     &A_r\le s\le A'.
\end{cases}
\]
The metric at the original scale is
\[
 \widetilde g=Q^{-1}\widehat g.
\]
\end{definition}

\begin{lemma}[Smoothness of the interpolating surgery metric]
\label{lem:interpolating-surgery-metric-smooth}
\dcref{MT:ch13:13.2}
\notready
The piecewise tensor $\widehat g$ of
\ref{def:interpolating-surgery-metric} is a smooth Riemannian metric on
$\mathcal S$.  It equals $Q\rho^*g$ on $s\le0$.
\end{lemma}

\begin{proof}
\uses{def:interpolating-surgery-metric}
The function $e^{-q/s}$, extended by zero for $s\le0$, is smooth and flat
at $s=0$.  Near $s=1$ and $s=2$, the cutoff $\alpha$ is constant with the
value required by the adjacent formula, and the same holds for $\beta$
near $s=A_r$ and $s=A'$.  Each middle tensor is a positive linear
combination of positive-definite metrics.  Near the tip, the last formula
is the constant multiple $e^{-2f(A')}\eta g_0$, so the nonsmoothness of the
radial coordinate at $p_0$ causes no loss of smoothness.  Finally $f=0$ on
$s\le0$, which gives the asserted agreement with the normalized neck.
\end{proof}

\subsection{Distance, curvature, and convergence}

\begin{lemma}[Nonexpansion of the standard radial collapse]
\label{lem:standard-cap-radial-collapse-nonexpanding}
\dcref{MT:ch12:12.37,MT:ch13:13.1}
\notready
In the data of \ref{def:local-neck-surgery-data}, define
$\xi\colon N\to\mathbb R^3$ by preserving the angular coordinate,
sending $s<A_0+4$ to the point at $g_0$-distance $A_0+4-s$ from $p_0$,
and collapsing $s\ge A_0+4$ to $p_0$.  For all sufficiently small
$\delta$,
\[
 \xi\colon(N,Qg)\longrightarrow(\mathbb R^3,\eta g_0)
\]
is distance nonincreasing.
\end{lemma}

\begin{proof}
\uses{def:local-neck-surgery-data}
The $\delta$-neck estimate and $\eta=1-2\delta$ give
\[
 Q\rho^*g\ge\eta(h_0+ds^2)
\]
as quadratic forms.  Write $g_0=dr^2+u(r)^2h_0$ in radial coordinates.
The standard cap has $|dr/ds|\le1$ and $u(r)\le1$.  Thus, for
$v=a\partial_s+w$ with $w\perp\partial_s$,
\[
 |d\xi(v)|_{g_0}^2
 =a^2|dr/ds|^2+u(r)^2|w|_{h_0}^2
 \le a^2+|w|_{h_0}^2.
\]
On the collapsed region $d\xi=0$.  Hence
$|d\xi(v)|_{\eta g_0}\le|v|_{Qg}$ almost everywhere.  Integrating along
piecewise $C^1$ curves and taking infima proves the claim.
\end{proof}

\begin{lemma}[The radial surgery map is nonexpanding]
\label{lem:radial-surgery-map-nonexpanding}
\dcref{MT:ch13:13.2}
\notready
Regard the map $\xi$ of
\ref{lem:standard-cap-radial-collapse-nonexpanding} as a map from $N$
to the surgery manifold $\mathcal S$.  Then
\[
 \xi\colon(N,g)\longrightarrow(\mathcal S,\widetilde g)
\]
is distance nonincreasing.
\end{lemma}

\begin{proof}
\uses{def:interpolating-surgery-metric,lem:standard-cap-radial-collapse-nonexpanding}
On $s\le0$ the normalized source and target metrics agree.  On
$0\le s\le1$, the target metric is $e^{-2f}Q\rho^*g\le Q\rho^*g$.
On $1\le s\le2$, the neck estimate gives
$\eta g_0\le Q\rho^*g$, so the convex interpolation, followed by the
factor $e^{-2f}$, is again no larger than $Q\rho^*g$.  On the remaining
cap, $\widehat g\le\eta g_0$, and
\ref{lem:standard-cap-radial-collapse-nonexpanding} applies.  Therefore
$\xi\colon(N,Qg)\to(\mathcal S,\widehat g)$ is nonexpanding.  Multiplying
both metrics by $Q^{-1}$ gives the statement.
\end{proof}

\begin{lemma}[Conformal curvature comparison on the neck side]
\label{lem:conformal-surgery-curvature-comparison}
\dcref{MT:ch13:13.5,MT:ch13:13.6,MT:ch13:13.7,MT:ch13:13.8,MT:ch13:13.9,MT:ch13:13.10,MT:ch13:13.11,MT:ch13:13.12}
\notready
Let $h$ be a metric $C^{\lfloor1/\delta\rfloor}$-close to
$h_0+ds^2$ on $S^2\times I$, and put $\widehat h=e^{-2f}h$ with
$f=C_0\delta e^{-q/s}$ for $s>0$.  Choose $K$ so that every curvature
eigenvalue of a scale-one $\delta$-neck is at least $-K\delta$.
There are universal choices $q\gg(A_0+4)^2$, $C_0=2Ke^q$, and
$\delta_*>0$ such that, whenever $0<\delta\le\delta_*$ and $0<s<4$,
the scalar curvature of
$\widehat h$ is at least that of $h$, and the least eigenvalue of its
curvature operator is strictly larger than that of $h$.
\end{lemma}

\begin{proof}
In coordinates with $x^0=s$, the derivatives of $f$ satisfy
\begin{align*}
 f_0&=\frac q{s^2}f,\\
 f_{00}&=\left(\frac{q^2}{s^4}-\frac{2q}{s^3}\right)f
          +\frac q{s^2}f\,O(\delta),\\
 f_{ij}&=\frac q{s^2}f\,O(\delta)\quad(i,j)\ne(0,0).
\end{align*}
Substitution in the conformal curvature formula gives
\[
 R_{\widehat h}=e^{2f}\left(
 R_h+4\left(\frac{q^2}{s^4}-\frac{2q}{s^3}\right)f
 -2\frac{q^2}{s^4}f^2+\frac q{s^2}f\,O(\delta)\right).
\]
Choose $q$ so that $q^2/s^4$ dominates $q/s^3$ on
$(0,A_0+4]$ while $q^2s^{-4}e^{-q/s}$ remains uniformly small.  Then
choose $\delta_*$ so that $f$, the $O(\delta)$ term, and the quadratic
term in $f$ are dominated by the positive linear term.  This proves
$R_{\widehat h}\ge R_h$.

After Gram--Schmidt and diagonalization of the small curvature block, let
$\lambda\ge\mu\ge\nu$ be the eigenvalues of $h$.  Neck closeness gives
$|\lambda-1/2|+|\mu|+|\nu|\le K\delta$.  Symmetric-matrix perturbation
applied to the same conformal formula yields
\begin{align*}
 \left|\lambda'-e^{2f}\left(\lambda-\frac{q^2}{s^4}f^2\right)\right|
   &\le A\frac{q^2}{s^4}f\delta,\\
 \left|\mu'-e^{2f}\left(\mu+\left(\frac{q^2}{s^4}-\frac{2q}{s^3}\right)f\right)\right|
   &\le A\frac{q^2}{s^4}f\delta,\\
 \left|\nu'-e^{2f}\left(\nu+\left(\frac{q^2}{s^4}-\frac{2q}{s^3}\right)f\right)\right|
   &\le A\frac{q^2}{s^4}f\delta.
\end{align*}
For the choices above,
\[
 \mu'\ge e^{2f}\left(\mu+\frac{q^2}{2s^4}f\right),
 \qquad
 \nu'\ge e^{2f}\left(\nu+\frac{q^2}{2s^4}f\right).
\]
The largest eigenvalue remains $1/2+o(1)$, while both small eigenvalues
increase strictly.  Hence the least eigenvalue increases.
\end{proof}

\begin{lemma}[Curvature and pinching on the retained neck]
\label{lem:retained-neck-surgery-pinching}
\dcref{MT:ch13:13.13,MT:ch13:13.14}
\notready
For the universal choices in
\ref{lem:conformal-surgery-curvature-comparison} and all sufficiently
small $\delta$, the metric $\widetilde g$ has positive sectional curvature
on $1\le s<4$.  Moreover, if the original neck satisfies, for some
$t\ge0$ and $X_g=\max(0,-\nu_g)$,
\[
 R_g\ge-\frac6{1+4t},
 \qquad
 R_g\ge2X_g\bigl(\log X_g+\log(1+t)-3\bigr)
 \quad(X_g>0),
\]
then the same inequalities hold for $\widetilde g$ on $s<4$.
\end{lemma}

\begin{proof}
\uses{def:interpolating-surgery-metric,lem:conformal-surgery-curvature-comparison}
On $1\le s<4$, the function
\[
 \frac{q^2f}{2s^4}
\]
is increasing, and at $s=1$ it equals
$q^2C_0e^{-q}\delta/2=q^2K\delta>K\delta$.  The eigenvalue estimates in
\ref{lem:conformal-surgery-curvature-comparison} therefore make both
small eigenvalues positive; the largest is greater than $1/4$ for small
$\delta$.

On $s\le0$ the metric is unchanged.  On $0<s<1$, the base metric in
\ref{def:interpolating-surgery-metric} is $Q\rho^*g$, so
\ref{lem:conformal-surgery-curvature-comparison}, followed by rescaling,
gives $R_{\widetilde g}\ge R_g$ and
$X_{\widetilde g}\le X_g$ at corresponding points.  If
$X_g\ge e^3/(1+t)$, the function
$u\mapsto2u(\log u+\log(1+t)-3)$ is increasing, so both pinching
inequalities pass to the new metric.  If $X_g<e^3/(1+t)$, the same is true
of $X_{\widetilde g}$ and the right-hand side is negative, whereas a
sufficiently round neck has nonnegative scalar curvature.  This proves the
claim on $s<1$.  On $1\le s<4$, the already established positive
sectional curvature gives $R_{\widetilde g}>0$ and
$X_{\widetilde g}=0$, so both pinching inequalities hold there as well.
\end{proof}

\begin{lemma}[Curvature and pinching on the standard-cap side]
\label{lem:standard-cap-surgery-curvature}
\dcref{MT:ch13:13.2}
\notready
For the same universal choices and all sufficiently small $\delta$, the
metric $\widetilde g$ has positive sectional curvature on
$4\le s\le A_0+4$.  It satisfies the two curvature-pinching inequalities
in \ref{lem:retained-neck-surgery-pinching} there.
\end{lemma}

\begin{proof}
\uses{def:standard-surgery-cap-model,def:interpolating-surgery-metric,lem:retained-neck-surgery-pinching}
On $4\le s\le A_0+4-r_0$, the metric is a conformal multiple of
$\eta g_0$.  The standard metric has nonnegative curvature.  In its radial
warped-product coordinates, the conformal curvature formula has positive
leading terms $q^2f/s^4$ in both the radial and tangential sectional
curvatures.  Since $g_0$ is fixed, choosing $q$ after $g_0$ makes these
terms dominate all bounded lower-order coefficients; decreasing $\delta$
then absorbs the quadratic terms in $f$.  On
$A_0+4-r_0\le s\le A_0+4$, the metric converges smoothly as
$\delta\to0$ to $g_0$, which has constant sectional curvature $1/4$ on
this region.  Positivity is open in the $C^2$ topology, so it persists for
small $\delta$.  Positive sectional curvature gives $X=0$ and positive
scalar curvature, which implies both pinching inequalities.
\end{proof}

\begin{lemma}[Convergence of the surgery cap to the standard cap]
\label{lem:local-surgery-standard-cap-convergence}
\dcref{MT:ch13:13.2}
\notready
For every $\delta''>0$ there is $\delta_1'(\delta'')>0$ such that, if
$\delta\le\delta_1'(\delta'')$, then $\widehat g$ on
$B_{\widehat g}(p_0,(\delta'')^{-1})$ is $\delta''$-close in
$C^{\lfloor1/\delta''\rfloor}$ to $g_0$ on the corresponding standard
ball.
\end{lemma}

\begin{proof}
\uses{def:interpolating-surgery-metric}
Fix $k=\lfloor1/\delta''\rfloor$ and the indicated standard ball.  The
cutoffs have fixed derivatives, $\eta\to1$, and every derivative through
order $k$ of $f$ tends uniformly to zero as $\delta\to0$.  On the portion
coming from the neck, the normalized metric $Q\rho^*g$ is
$\delta$-close to the product cylinder, which is exactly $g_0$ outside the
standard cap.  Thus every piece of $\widehat g$ converges to $g_0$ in
$C^k$, uniformly across the cutoff regions.  Smooth convergence also makes
the two metric balls correspond after decreasing $\delta$, proving the
stated closeness.
\end{proof}

\begin{theorem}[Local surgery on a delta-neck]
\label{thm:local-neck-surgery}
\dcref{MT:ch13:13.2}
\notready
There are universal constants $C_0,q,R_0<\infty$ and $\delta_0'>0$ such
that, whenever $Q=R_g(x_0)\ge R_0$ and
$0<\delta\le\delta_0'$, the metric $\widetilde g$ of
\ref{def:interpolating-surgery-metric} has the following properties.
\begin{enumerate}
\item Every curvature-pinching inequality in
      \ref{lem:retained-neck-surgery-pinching} that holds on the original
      neck holds everywhere on $\mathcal S$.
\item The sectional curvature is positive on
      $1\le s\le A_0+4$.
\item The radial collapse map
      $\xi\colon(N,g)\to(\mathcal S,\widetilde g)$ is distance
      nonincreasing.
\item The normalized pointed cap converges smoothly on bounded balls,
      \[
        (\mathcal S,\widehat g,p_0)\longrightarrow
        (\mathbb R^3,g_0,p_0),
      \]
      as $\delta\to0$, in the quantitative sense of
      \ref{lem:local-surgery-standard-cap-convergence}.
\end{enumerate}
\end{theorem}

\begin{proof}
\uses{def:interpolating-surgery-metric,lem:interpolating-surgery-metric-smooth,lem:radial-surgery-map-nonexpanding,lem:retained-neck-surgery-pinching,lem:standard-cap-surgery-curvature,lem:local-surgery-standard-cap-convergence}
Choose the common $q$ and $C_0$ required by the curvature lemmas, and then
choose $\delta_0'$ small enough for all cited lemmas.  The two curvature
lemmas cover respectively $s<4$ and $s\ge4$, proving the first two items.
The third item is \ref{lem:radial-surgery-map-nonexpanding}, and the fourth
is \ref{lem:local-surgery-standard-cap-convergence}.  Smoothness is supplied
by \ref{lem:interpolating-surgery-metric-smooth}.
\end{proof}

\section{Canonical Neck and Cap Topology}

\subsection{Canonical regions and neck overlap}

\begin{definition}[Canonical cap]
\label{def:canonical-cap}
\dcref{MT:ch9:9.72}
\uses{def:epsilon-neck-and-scale}
\notready
For $C<\infty$, a \emph{$(C,\epsilon)$-cap} is an open submanifold
$\mathcal C$ of a Riemannian three-manifold with an open
$\epsilon$-neck $N\subset\mathcal C$ such that:
\begin{enumerate}
\item $\mathcal C$ is diffeomorphic to an open three-ball or to
      $\mathbb{RP}^3$ minus a point;
\item $\overline Y=\mathcal C\setminus N$ is a compact submanifold with
      boundary, its interior $Y$ is the \emph{core}, and
      $\partial\overline Y$ is a central sphere of an $\epsilon$-neck;
\item $R>0$ on $\mathcal C$,
      $\sup_{x,y\in\mathcal C}R(y)/R(x)<C$, and
      $\operatorname{diam}(\mathcal C)
       <C(\sup_{\mathcal C}R)^{-1/2}$.
\end{enumerate}
\end{definition}

\begin{definition}[Tubes, capped tubes, and neck fibrations]
\label{def:canonical-neck-cap-regions}
\dcref{MT:ch9:9.81,MT:app:A.21,MT:app:A.25}
\uses{def:epsilon-neck-and-scale,def:canonical-cap}
\notready
An \emph{$\epsilon$-tube} is a submanifold diffeomorphic to
$S^2\times I$, for a nondegenerate interval $I$, covered by
$\epsilon$-necks whose central spheres are isotopic to the $S^2$ factors.
An open tube has no boundary.  A \emph{$C$-capped $\epsilon$-tube} is the
connected union of an open $\epsilon$-tube and one $(C,\epsilon)$-cap,
glued along an end of each.  A \emph{doubly $C$-capped
$\epsilon$-tube} has disjoint cap cores attached at both ends.  An
\emph{$\epsilon$-fibration} is a closed connected $S^2$-bundle over
$S^1$, covered by $\epsilon$-necks whose central spheres are isotopic to
the fibers.
\end{definition}

\begin{definition}[Balanced chain of epsilon-necks]
\label{def:balanced-neck-chain}
\dcref{MT:app:A.12,MT:app:A.17}
\uses{def:epsilon-neck-and-scale}
\notready
A finite \emph{chain of $\epsilon$-necks} is a sequence
$N_a,\ldots,N_b$ with oriented axial coordinates such that
$N_i\cap N_{i+1}$ contains both
\[
 \{s_{N_i}>\tfrac12\epsilon^{-1}\}
 \quad\hbox{and}\quad
 \{s_{N_{i+1}}<-\tfrac12\epsilon^{-1}\},
\]
is contained in
\[
 \{s_{N_i}>-\tfrac12\epsilon^{-1}\}
 \cap
 \{s_{N_{i+1}}<\tfrac12\epsilon^{-1}\},
\]
and $N_i\cap E_{N_a}^-=\varnothing$ for $i>a$.  An infinite chain indexed
by an interval in $\mathbb Z$ is one whose restriction to every finite
subinterval is a chain.  The chain is \emph{balanced} if consecutive
centers $x_j,x_{j+1}$ satisfy
\[
 0.99\,r_{N_j}\epsilon^{-1}
 \le d(x_j,x_{j+1})
 \le1.01\,r_{N_j}\epsilon^{-1}.
\]
\end{definition}

\begin{lemma}[Metric control in a small epsilon-neck]
\label{lem:epsilon-neck-metric-control}
\dcref{MT:app:A.2,MT:app:A.5}
\notready
Write $d_N$ for the intrinsic distance induced by the restriction of the
ambient metric to $N$.  The estimates below are intrinsic to the neck; no
ambient no-shortcut assertion is being made.
For every $0<\alpha<1/8$, all sufficiently small $\epsilon$ have the
following property.  If $N$ is an $\epsilon$-neck of scale $r_N$ centered
at $x$, then scalar curvatures on $N$ lie between
$(1-\alpha)r_N^{-2}$ and $(1+\alpha)r_N^{-2}$.  Neck spheres have
diameter at most $(1+\alpha)\sqrt2\pi r_N$, and points with sufficiently
large axial separation, namely
$|s_N(q)-s_N(p)|\ge\epsilon^{-1}/100$ or
$d_N(p,q)\ge r_N\epsilon^{-1}/100$, satisfy
\[
 (1-\alpha)r_N|s_N(q)-s_N(p)|
 \le d_N(p,q)
 \le(1+\alpha)r_N|s_N(q)-s_N(p)|.
\]
In particular, every path contained in $N$ and joining its two end regions has length at least
$2(1-\alpha)r_N\epsilon^{-1}$.
\end{lemma}

\begin{proof}
\uses{def:epsilon-neck-and-scale}
After multiplying the metric by $r_N^{-2}$, the neck is arbitrarily close
in $C^2$ to the round product cylinder.  Curvature continuity gives the
scalar-curvature bounds and the diameter bound for the slices.  In the
product cylinder, a minimizing geodesic has an axial projection of speed at
most one, while its spherical projection has uniformly bounded diameter.
For an axial displacement tending to infinity with $\epsilon^{-1}$, the
spherical error is absorbed into the factor $1\pm\alpha$.  The same
argument applies to intrinsic distance in $N$ when a large intrinsic
distance, rather than a large axial displacement, is assumed, since the
spherical diameter is bounded.
Rescaling by $r_N$ proves the distance estimate.  A path exiting both ends
has axial displacement $2\epsilon^{-1}$, which gives the last assertion.
\end{proof}

\begin{proposition}[Overlap control for small epsilon-necks]
\label{prop:overlapping-epsilon-necks}
\dcref{MT:app:A.11}
\notready
For every $0<\alpha\le10^{-2}$, all sufficiently small $\epsilon$ have
the following property.  If two $\epsilon$-necks $N,N'$ intersect in a region
meeting the middle nine-tenths of each neck, then
their scales differ by a factor in $(1-\alpha,1+\alpha)$, their neck-sphere
tangent planes make angle less than $\alpha$ on the overlap, and their
oriented axial directions agree up to sign and an error less than
$\alpha$.  A neck sphere of one neck contained in the other is isotopic
there to a neck sphere.  If the overlap meets the middle nine-tenths of one
neck, then $N\cap N'$ contains neck spheres from both structures which are
isotopic within the overlap, and the relevant component of the overlap is
diffeomorphic to $S^2\times(0,1)$.
\end{proposition}

\begin{proof}
\uses{def:epsilon-neck-and-scale,lem:epsilon-neck-metric-control}
At an intersection point, both normalized metrics are close to the same
round-cylinder metric.  Comparing scalar curvature there gives the scale
ratio.  On a sufficiently round cylinder the curvature operator has a
simple eigenvalue near $1/2$ whose eigenspace is the tangent plane of the
sphere factor; continuity of this eigenspace proves the tangent-plane
estimate.  Its orthogonal line gives the axial directions, determined up
to sign.

If a sphere $S'$ from $N'$ lies in $N$, projection along the axial lines of
$N$ restricts to a local diffeomorphism $S'\to S^2$.  It is a covering
because $S'$ is compact, and hence a diffeomorphism because $S^2$ is simply
connected.  Thus $S'$ is a graph over a neck sphere and is isotopic to it
along axial lines.  In a middle overlap, the metric and axial-coordinate
estimates provide such graph spheres from both neck structures.  The
isotopy between them, together with their collars, supplies the asserted
product structure on the overlap.
\end{proof}

\subsection{Chains of necks}

\begin{lemma}[A neck chain is a product]
\label{lem:neck-chain-product}
\dcref{MT:app:A.13}
\notready
The union $U$ of a finite or infinite chain of sufficiently small
$\epsilon$-necks is diffeomorphic to $S^2\times(0,1)$, compatibly with
the isotopy classes of all neck spheres.
\end{lemma}

\begin{proof}
\uses{def:balanced-neck-chain,prop:overlapping-epsilon-necks}
For a finite chain, induct on the number of necks.  When the last neck is
adjoined, \ref{prop:overlapping-epsilon-necks} identifies the overlap with
$S^2\times(0,1)$ and identifies a cross-sectional sphere with the factor
sphere in both pieces.  Gluing the corresponding collars therefore extends
the product coordinate across the new neck.  The construction can be made
compatible under inclusion.  Exhausting an infinite chain by finite
subchains then gives a proper interval coordinate on the union and hence a
diffeomorphism $U\cong S^2\times(0,1)$.
\end{proof}

\begin{lemma}[An infinite neck-chain end has no frontier]
\label{lem:infinite-neck-chain-frontier}
\dcref{MT:app:A.15}
\notready
Let $N_0,N_1,\ldots$ be a chain of sufficiently small
$\epsilon$-necks in $M$, and put $U=\bigcup_{j\ge0}N_j$.  Then the
frontier of $U$ in $M$ is exactly the frontier of the negative end of
$N_0$.
\end{lemma}

\begin{proof}
\uses{def:balanced-neck-chain,lem:epsilon-neck-metric-control}
The frontier of a finite subchain consists of the two exposed end
frontiers.  Hence a limit point of $U$ represented by points in one fixed
finite subchain can lie in the frontier of $U$ only at the negative end of
$N_0$: the positive end of every finite subchain lies in the next neck.

It remains to exclude a convergent sequence $z_k\in N_{j_k}$ with
$j_k\to\infty$.  If $z_k\to z$, continuity of scalar curvature bounds
$R(z_k)$ and gives a positive lower bound for the scales of all sufficiently
late necks.  By \ref{lem:epsilon-neck-metric-control}, those necks have a
uniformly positive width.  Any path from $z_k$ to $z$ must cross either
$N_0$ or a complete sufficiently late neck, so $d(z_k,z)$ has a fixed
positive lower bound, a contradiction.  Thus the infinite end contributes
no frontier.
\end{proof}

\begin{lemma}[Extension of a balanced neck chain]
\label{lem:balanced-neck-chain-extension}
\dcref{MT:app:A.17,MT:app:A.18}
\notready
For all sufficiently small $\epsilon$, let
$N_a,\ldots,N_b$ be a balanced chain whose central spheres separate $M$.
If a point in the frontier of its positive end is the center of an
$\epsilon$-neck $N$, then, after reversing $N$ if necessary,
$N_a,\ldots,N_b,N$ is a balanced chain.  The analogous assertion holds at
the negative end.
\end{lemma}

\begin{proof}
\uses{def:balanced-neck-chain,lem:epsilon-neck-metric-control,prop:overlapping-epsilon-necks}
The two cases are symmetric.  At the positive end, the new center lies in
the closure of $N_b$.  The distance estimate in
\ref{lem:epsilon-neck-metric-control} gives the balanced inequalities for
the two centers, and \ref{prop:overlapping-epsilon-necks} allows the axial
coordinate of $N$ to be oriented consistently with $N_b$.  The separating
central sphere of $N_a$ prevents $N$ from returning to the negative end of
the chain: otherwise a path through the overlaps would cross that sphere
exactly once.  Thus all defining conditions of a balanced chain remain
valid.
\end{proof}

\begin{theorem}[Separating neck centers lie in an epsilon-tube]
\label{thm:separating-neck-centers-form-tube}
\dcref{MT:app:A.19}
\notready
There is $\epsilon_0>0$ such that the following holds for
$0<\epsilon\le\epsilon_0$.  Let $X$ be a connected subset of a Riemannian
three-manifold $M$.  Suppose every $x\in X$ is the center of an
$\epsilon$-neck and every corresponding central sphere separates $M$.
Then a subcollection of these necks, after reversing axial coordinates when
necessary, is a balanced chain whose union contains $X$.  Consequently
$X$ lies in an $\epsilon$-tube.
\end{theorem}

\begin{proof}
\uses{def:canonical-neck-cap-regions,lem:neck-chain-product,lem:infinite-neck-chain-frontier,lem:balanced-neck-chain-extension}
Start with one neck centered in $X$.  If a finite balanced chain does not
contain $X$, connectedness supplies a point of $X$ on the frontier of one
of its two ends.  By \ref{lem:balanced-neck-chain-extension}, a neck
centered there extends the chain.  Repeat at both ends.

If the process terminates, the union contains $X$.  For a half-infinite
chain, \ref{lem:infinite-neck-chain-frontier} says that only its finite end
has frontier; if this end cannot be extended, connectedness again forces
$X$ into the union.  A chain infinite in both directions has no frontier
and its union is an ambient component.  In every case
\ref{lem:neck-chain-product} identifies the union with
$S^2\times(0,1)$, so it is an $\epsilon$-tube.
\end{proof}

\subsection{Caps and classification}

\begin{lemma}[No infinite nested chain of canonical caps]
\label{lem:no-infinite-nested-canonical-caps}
\dcref{MT:app:A.22}
\notready
For fixed $C$ and sufficiently small $\epsilon$, there is no chain of
$(C,\epsilon)$-caps
\[
 C_0\subset C_1\subset\cdots
\]
such that the closure of the core of $C_i$ meets
$\operatorname{Fr}_M(C_{i-1})$ for every $i\ge1$.
\end{lemma}

\begin{proof}
\uses{def:canonical-cap,lem:epsilon-neck-metric-control,prop:overlapping-epsilon-necks}
Fix $x_0\in C_0$ and put $Q_0=R(x_0)$.  Let $S_i'$ be the boundary of
the core of $C_i$, realized as the central sphere of a neck $N_i'$, and
let $N_i$ be the neck end of $C_i$.  Curvature comparability in a cap
gives every $N_i$ scale at least
$C^{-1/2}Q_0^{-1/2}$.

If $S_i'$ is disjoint from $C_{i-1}$, the component of
$M\setminus S_i'$ containing $C_{i-1}$ must be the core side; otherwise
$C_{i-1}\cap C_i$ would lie entirely in the neck end, contrary to
$C_{i-1}\subset C_i$.  Hence the core of $C_i$ contains $C_{i-1}$, and
\ref{lem:epsilon-neck-metric-control} gives
\[
 d(x_0,M\setminus C_i)-d(x_0,M\setminus C_{i-1})
 \ge c\epsilon^{-1}C^{-1/2}Q_0^{-1/2}
\]
for a universal $c>0$.  If $S_i'\subset C_{i-1}$, the component of its
complement whose closure lies in $C_{i-1}$ cannot be the core side, since
that core meets $\operatorname{Fr}(C_{i-1})$.  It would therefore contain
the neck end of $C_i$ while having only one boundary sphere, contradicting
the two frontier components of a neck end.

In the remaining case $S_i'$ crosses $N_{i-1}$.
By \ref{prop:overlapping-epsilon-necks}, the cylinder directions almost
agree or almost oppose.  Separation forces the agreeing orientation, so
the central sphere of $N_i$ lies beyond the positive end of $N_{i-1}$.
The same neck-width estimate gives the displayed increase.  Thus every
noncontradictory step raises the distance to the complement by a fixed
positive amount.  But \ref{def:canonical-cap} bounds
$\operatorname{diam}(C_i)$ by $CQ_0^{-1/2}$.  This is impossible for
arbitrarily large $i$.
\end{proof}

\begin{lemma}[An isotopic cap-gluing sphere determines a component]
\label{lem:isotopic-cap-gluing-sphere}
\dcref{MT:app:A.24}
\notready
Let $X\subset M$ be connected.  Let $\widetilde C$ be a cap or a finite
capped tube containing a point of $X$, and let $C'$ be a
$(C,\epsilon)$-cap whose core contains a point of
$X\cap\operatorname{Fr}_M(\widetilde C)$.  Let $N'$ be the neck end of
$C'$, oriented away from its core.  If a two-sphere
$\Sigma\subset N'\cap\widetilde C$ lies in the closed positive half of
$N'$ and is isotopic in $N'$ to its central sphere, then
$\widetilde C\cup C'$ is a component of $M$ containing $X$.
\end{lemma}

\begin{proof}
\uses{def:canonical-cap,def:canonical-neck-cap-regions}
The sphere $\Sigma$ separates $\widetilde C$ into a relatively compact
side $A$ and an end side.  It also separates $C'$; let $A'$ be the side
with compact closure, which contains the core of $C'$ because $\Sigma$ is
isotopic to the central sphere of $N'$.  Both $A$ and $A'$ have frontier
$\Sigma$.  They cannot lie on the same side of $\Sigma$, since then the
core of $C'$ would lie in $\widetilde C$, contrary to its meeting the
frontier of $\widetilde C$.  Thus they lie on opposite sides, and their
closures form an ambient component.  This component is
$\widetilde C\cup C'$.  It meets $X$, so connectedness of $X$ places all
of $X$ in it.
\end{proof}

\begin{lemma}[Adjacent canonical caps form an ambient component]
\label{lem:adjacent-canonical-caps-form-component}
\dcref{MT:app:A.23}
\notready
Let $X\subset M$ be connected, and let $\widetilde C$ be obtained from a
$(C,\epsilon)$-cap $C_0$ by adjoining a finite balanced chain to its neck
end.  Suppose $\widetilde C$ contains a point of $X$.  Let $C'$ be a
$(C,\epsilon)$-cap whose core contains a point of
$X\cap\operatorname{Fr}_M(\widetilde C)$, and assume
$C_0\nsubseteq C'$.  Then $\widetilde C\cup C'$ is a component of $M$
containing $X$.
\end{lemma}

\begin{proof}
\uses{def:canonical-cap,def:balanced-neck-chain,prop:overlapping-epsilon-necks,lem:isotopic-cap-gluing-sphere}
Let $S'$ be the boundary of the core of $C'$ and $N'$ its neck end.
If $S'\subset\widetilde C$, apply
\ref{lem:isotopic-cap-gluing-sphere} with $\Sigma=S'$.  If $S'$ is
disjoint from $\widetilde C$, the side of $S'$ containing
$\widetilde C$ cannot be the core side, since that would imply
$C_0\subset C'$.  Hence the neck $N'$ meets the last neck of the balanced
chain.  By \ref{prop:overlapping-epsilon-necks}, their overlap contains a
sphere $\Sigma$ isotopic to $S'$ and lying in the positive half of $N'$.
Apply \ref{lem:isotopic-cap-gluing-sphere}.

Finally suppose $S'$ crosses the last neck of $\widetilde C$.  The scale,
diameter, and orientation conclusions of
\ref{prop:overlapping-epsilon-necks} again give a common sphere
$\Sigma$ in the positive half of $N'$.  If the cap core and
$\widetilde C$ lie on opposite sides of $\Sigma$, apply
\ref{lem:isotopic-cap-gluing-sphere}.  If they lie on the same side, the
core of $C'$ contains $\widetilde C$ up to its last half-neck; the overlap
then supplies that half-neck as well, forcing $C_0\subset C'$, a
contradiction.  These cases are exhaustive.
\end{proof}

\begin{theorem}[Connected sets covered by canonical cap cores and neck centers]
\label{thm:cap-and-neck-cover-classification}
\dcref{MT:app:A.21}
\notready
For every $C<\infty$ and all sufficiently small $\epsilon>0$, let $X$ be
a connected subset of a Riemannian three-manifold $M$.  Suppose every
point of $X$ is either the center of an $\epsilon$-neck or lies in the core
of a $(C,\epsilon)$-cap.  Then at least one of the following holds:
\begin{enumerate}
\item $X$ lies in a component which is the union of two
      $(C,\epsilon)$-caps and is diffeomorphic to
      $S^3$, $\mathbb{RP}^3$, or
      $\mathbb{RP}^3\#\mathbb{RP}^3$;
\item $X$ lies in a doubly $C$-capped $\epsilon$-tube component of one
      of those three diffeomorphism types;
\item $X$ lies in one $(C,\epsilon)$-cap;
\item $X$ lies in a $C$-capped $\epsilon$-tube;
\item $X$ lies in an $\epsilon$-tube;
\item $X$ lies in an $\epsilon$-fibration component covered by
      $\epsilon$-necks.
\end{enumerate}
\end{theorem}

\begin{proof}
\uses{def:canonical-neck-cap-regions,lem:neck-chain-product,lem:infinite-neck-chain-frontier,thm:separating-neck-centers-form-tube,lem:no-infinite-nested-canonical-caps,lem:adjacent-canonical-caps-form-component,prop:overlapping-epsilon-necks}
Suppose first that $x_0\in X$ lies in a cap core.  Starting with a cap
containing $x_0$, enlarge it whenever a larger cap has core meeting the
relevant frontier.  By
\ref{lem:no-infinite-nested-canonical-caps}, this process reaches a cap
$C_0$ for which no further enlargement is possible.

If $X$ misses the frontier of $C_0$, connectedness puts $X$ in $C_0$,
giving alternative (3).  If the frontier points of $X$ are neck centers,
adjoin a balanced chain.  A finite chain which stops, or a chain with an
infinite end, forms with $C_0$ a capped tube containing $X$ by
\ref{lem:infinite-neck-chain-frontier}; this is alternative (4).

The remaining possibility is that a cap $C'$ has core meeting the frontier
of the capped chain.  Maximality of $C_0$ gives $C_0\nsubseteq C'$, so
\ref{lem:adjacent-canonical-caps-form-component} shows that the union is
an entire component containing $X$.  The product structure from
\ref{lem:neck-chain-product} reduces its topology to gluing two cap cores
along a two-sphere.  Each core is a three-ball or a punctured
$\mathbb{RP}^3$, so the result is $S^3$, $\mathbb{RP}^3$, or
$\mathbb{RP}^3\#\mathbb{RP}^3$.  According as the intervening neck chain
is absorbed into the two cap descriptions or retained, this gives
alternative (1) or (2).

Now suppose every point of $X$ is a neck center.  If all central spheres
separate $M$, \ref{thm:separating-neck-centers-form-tube} gives
alternative (5).  Otherwise extend a balanced chain of necks centered on
$X$.  A chain which terminates or has an infinite product end contains
$X$, by \ref{lem:infinite-neck-chain-frontier}, and gives (5).  The only
remaining case is that a finite chain returns to its first end.
By \ref{prop:overlapping-epsilon-necks}, the returning overlap closes the
sphere product, directly or after lifting to the universal cover.  The deck
transformation preserves the isotopy classes of the cross-sectional
spheres, so the component is an $S^2$-bundle over $S^1$ covered by necks.
This is alternative (6).
\end{proof}

\begin{theorem}[Topology of a manifold covered by canonical caps and necks]
\label{thm:canonical-neck-cap-topology}
\dcref{MT:app:A.25}
\notready
Let $(M,g)$ be a connected Riemannian three-manifold.  For fixed
$C<\infty$ and sufficiently small $\epsilon>0$, suppose each point of
$M$ is either the center of an $\epsilon$-neck or belongs to the core of a
$(C,\epsilon)$-cap.  Exactly one of the following types occurs:
\begin{enumerate}
\item $M$ is diffeomorphic to $S^3$, $\mathbb{RP}^3$, or
      $\mathbb{RP}^3\#\mathbb{RP}^3$, and is the union of two caps or a
      doubly capped $\epsilon$-tube;
\item $M$ is diffeomorphic to $\mathbb R^3$ or to
      $\mathbb{RP}^3\setminus\{p\}$, and is a cap or a capped
      $\epsilon$-tube;
\item $M$ is diffeomorphic to $S^2\times\mathbb R$ and is an
      $\epsilon$-tube;
\item $M$ is diffeomorphic to an $S^2$-bundle over $S^1$ and is an
      $\epsilon$-fibration.
\end{enumerate}
\end{theorem}

\begin{proof}
\uses{thm:cap-and-neck-cover-classification}
Apply \ref{thm:cap-and-neck-cover-classification} with $X=M$.  Since $M$
is a whole connected component, alternatives (1) and (2) give the first
type.  Alternatives (3) and (4) give the second type, using the two possible
cap-core topologies.  Alternative (5), together with the product definition
of a tube, gives the third type, and alternative (6) gives the fourth.
\end{proof}
